\documentclass[phd,english]{../uw-thesis2}

\title{Title of the doctoral dissertation}
\author{Your Full Name}
\date{March 2026}

\faculty{Faculty of Mathematics, Informatics, and Mechanics}
\field{Natural Sciences}
\fieldpl{nauk ścisłych i przyrodniczych}
\discipline{Computer and Information Sciences}
\disciplinepl{informatyka}
\supervisor{Prof.~Name Surname}
\cosupervisor{Dr~Name Surname} % remove this line if there is no assistant supervisor
\keywords{keyword one; keyword two; keyword three}
\polishtitle{Polski tytuł rozprawy doktorskiej}
\polishabstract{Tutaj wstaw streszczenie po polsku.}
\englishabstract{Put the English abstract here.}
\titlepagelogo{../UW_logo.pdf} % optional; omitted automatically if the file does not exist

\newtheorem{theorem}{Theorem}[chapter]
\newtheorem{lemma}[theorem]{Lemma}

\begin{document}
\makefrontmatter

\chapter{Introduction}
Start the thesis here.

\section{First section}

\begin{theorem}
This is a sample theorem. You can use the \texttt{theorem} environment for theorems, lemmas, propositions, etc.
\end{theorem}

\cite{tarski_decision}

\section{Second section}

\begin{theorem}
This is a second theorem!
\end{theorem}

\begin{lemma}
This is a sample lemma.
\end{lemma}

\chapter{Contents}

This is the second chapter.

\begin{lemma}
This is a sample lemma.
\end{lemma}

\begin{lemma}
This is a sample lemma.
\end{lemma}

\bibliographystyle{alpha}
\bibliography{test.bib}
\end{document}
